La observación de la superficie terrestre mediante imágenes satelitales es importante en múltiples tareas de monitoreo ambiental, análisis territorial, agricultura, etc.

Sin embargo, una de las principales limitaciones de las imágenes ópticas satelitales es la presencia de nubes. En muchas ocasiones, una parte significativa de las capturas puede encontrarse parcial o totalmente cubierta por nubosidad. Esto reduce la disponibilidad de imágenes útiles y dificulta el análisis temporal de una misma región. En este contexto, la tarea de \textit{cloud removal} busca reconstruir una imagen libre de nubes a partir de una observación degradada, intentando recuperar la mayor cantidad posible de información espacial y espectral de la escena original.

El problema presenta una dificultad particular, ya que la información oculta por las nubes no se encuentra directamente disponible en la imagen óptica de entrada. En este trabajo se utiliza información SAR proveniente de Sentinel-1, ya que este tipo de sensor opera mediante radar y puede capturar características de la superficie terrestre aun en presencia de nubes, sin depender de la iluminación solar ni verse afectado de la misma forma por la nubosidad.  Sin embargo, no siempre se cuenta con este tipo de capturas, por lo que resulta relevante estudiar en qué medida la información SAR aporta o no una mejora en la reconstrucción.


% Por este motivo, resulta necesario incorporar modelos capaces de inferir contenido faltante a partir del contexto visible y, cuando es posible, aprovechar fuentes complementarias de información. En este trabajo se utiliza información SAR proveniente de Sentinel-1, ya que este tipo de sensor opera mediante radar y puede capturar características de la superficie terrestre aun en presencia de nubes. A diferencia de las imágenes ópticas, las imágenes SAR no dependen de la iluminación solar ni se ven afectadas de la misma forma por la nubosidad, por lo que constituyen una modalidad especialmente útil para complementar la reconstrucción.

En los últimos años, los modelos difusivos mostraron resultados relevantes en tareas generativas y de restauración de imágenes. Su capacidad para modelar procesos progresivos de degradación y reconstrucción los vuelve una alternativa relevante para problemas donde se busca recuperar una señal limpia a partir de una observación alterada. No obstante, en el caso específico de remoción de nubes, partir desde ruido gaussiano puro puede resultar innecesariamente costoso, ya que la imagen nublada todavía conserva parte de la estructura de la escena. Por este motivo, enfoques como DB-CR proponen formular el problema como un puente de difusión entre la distribución de imágenes nubladas y la distribución de imágenes libres de nubes.

El objetivo de este trabajo es implementar y comparar distintas arquitecturas basadas en modelos difusivos para la remoción de nubes en imágenes satelitales. Para ello, se parte de un modelo DDPM condicional y luego se desarrollan variantes inspiradas en DBCR, evaluando tanto modelos sin información SAR como distintas estrategias de incorporación multimodal. En particular, se analizan tres formas de utilizar la señal radar: concatenación directa de canales, una rama auxiliar inspirada en ControlNet y una arquitectura más compleja basada en bloques de fusión SAR-óptica mediante atención cruzada.

A lo largo del informe se describe el conjunto de datos utilizado, el marco teórico asociado a los modelos difusivos y a la formulación de \textit{diffusion bridge}, las arquitecturas implementadas, las métricas de entrenamiento y evaluación y los resultados obtenidos en cada experimento. Finalmente, se comparan las distintas variantes propuestas con el fin de analizar si la información SAR mejora la reconstrucción óptica y qué tipo de fusión resulta más efectiva bajo las restricciones computacionales del trabajo.